\paragraph{Scope note.}
This paper (E.4) concerns \emph{within-district} partisan homogeneity under single-member districts: how to deliberately construct geographically compact safe seats as a legal baseline for courts.
A companion paper (E.5) addresses the orthogonal question of \emph{between-party} allocation: how party-specific overlapping districts achieve proportionality through D'Hondt seat allocation before districting begins.
E.4 uses standard single-member geography; E.5 uses non-geographic party-specific districts.
The two papers share the theme of partisan district engineering but address different structural questions and different audiences (courts vs.\ legislatures seeking proportional reform).

Safe congressional districts---districts where one party's presidential vote
share consistently exceeds 60\%---now account for approximately 90\% of all
House seats \citep{cook2020competitive}.  This homogeneity arises from two
sources: organic geographic sorting (Democrats concentrate in cities,
Republicans in rural areas) \citep{bishop2009bigSort}, and deliberate
partisan gerrymandering (mapmakers pack opposition voters and crack moderate
constituencies) \citep{stephanopoulos2015partisan}.  The two mechanisms are
observationally difficult to separate: a state with many safe districts might
reflect geography, intentional manipulation, or a combination of both.

\subsection{The Safe Seat Debate}

There are two competing normative visions of redistricting.  The
\emph{competitive districts} view holds that elections should be contested,
that safe seats breed complacency and extremism, and that democratic
accountability requires competitive races \citep{jacobson2015politics}.  The
\emph{representation} view holds that safe districts ensure minority
representation (majority-minority districts are deliberately safe for
minority-preferred candidates) and that forcing artificial competition
misrepresents naturally sorted communities \citep{rush2000redistricting}.

This paper takes no position in this normative debate.  Instead, we contribute
a methodological tool: an algorithm that \emph{deliberately} generates safe
seats in a transparent, reproducible, and legally auditable way.  The value
of this tool is not to advocate for safe seats but to establish a neutral
reference point---a ``maximum-safe-seat'' baseline---against which enacted
plans can be compared.

\subsection{The Inversion Idea}

The standard bisect algorithm minimises edge cuts, which corresponds to
minimising district perimeter and maximising geographic compactness.  In
this mode, neighbouring tracts are assigned to the same district because they
share a long boundary, not because they have similar partisan lean.

We propose a simple inversion: modify the edge weights to reflect
\emph{partisan similarity} rather than shared boundary length.  Specifically,
we set:
\begin{equation}
  w_{uv} = \exp\!\bigl(-\alpha \cdot |D_u - D_v|\bigr)
  \label{eq:psd-weight}
\end{equation}
where $D_u$ and $D_v$ are the Democratic presidential vote shares in tracts
$u$ and $v$, and $\alpha > 0$ is a scaling parameter.  When $\alpha = 0$,
equation~\eqref{eq:psd-weight} reduces to unit weights (unweighted partitioning).
As $\alpha \to \infty$, only edges between tracts with identical partisan
lean receive nonzero weight, producing maximally homogeneous districts.

This ``partisan similarity'' weighting is the \emph{inverse} of what bisect
normally does: instead of attracting geographically proximate tracts, we
attract partisan-similar tracts.  The result is partisan similarity districts
(PSDs): districts that cluster geographically contiguous tracts with similar
partisan lean, producing homogeneous safe seats.

\subsection{Why PSDs Matter}

PSDs have three important uses:

\paragraph{Legal baseline.}
\emph{Rucho v.\ Common Cause} (2019) held that federal courts cannot adjudicate
partisan gerrymandering claims under the Constitution, leaving state courts
and legislatures as the primary venues for such challenges.  A neutral
algorithmic method for generating maximally safe seats---one that distributes
safe seats \emph{proportionally} between parties---provides a natural baseline:
an enacted plan that exceeds the PSD safe-seat count for one party, while
being below it for the other, provides quantitative evidence of intentional
imbalance.

\paragraph{Efficiency Gap calibration.}
The Efficiency Gap (EG) measures partisan asymmetry in wasted votes
\citep{stephanopoulos2015partisan}.  PSDs allow analysts to compute the EG
that would result from a ``symmetric'' safe-seat map, establishing the
maximum EG achievable without intentional partisan skew.

\paragraph{Algorithm benchmarking.}
PSDs provide a worst-case compactness benchmark.  Any redistricting algorithm
that produces districts less compact than PSDs is generating shapes not
justified by partisan-homogeneity goals; PSDs establish the compactness floor
for safe-seat-focused redistricting.

\subsection{Paper Structure}

Section~2 reviews the partisan clustering literature, the Efficiency Gap, and
geographic sorting.  Section~3 describes the PSD weighting formula and
implementation.  Section~4 presents results for North Carolina and Texas,
comparing PSDs to enacted plans and to neutral bisect plans.  Section~5
analyses the legal implications under \emph{Rucho} and related jurisprudence.
Section~6 concludes.
